% ============================================================================
% COURS 13 — MACHINES ASYNCHRONES ET SYNCHRONES
% Partie C : machine synchrone
% Source : 17-MAS_MS.docx
% ============================================================================

\section{Machine synchrone}
\label{sec:masms-5}

Une machine synchrone est un convertisseur électromécanique réversible. Elle fonctionne en moteur ou en génératrice, appelée alors \textbf{alternateur}. Le rotor crée un champ magnétique par aimants permanents ou par un enroulement d'excitation. En régime établi, il tourne exactement à la vitesse du champ tournant statorique :
\[
\boxed{\Omega=\Omega_s=\frac{\omega}{p}.}
\]

\subsection{Domaines d'emploi}

Les machines synchrones couvrent une très large gamme de puissances :
\begin{itemize}
  \item petites puissances : modélisme, ventilateurs, dispositifs médicaux et micromécanismes ;
  \item puissances moyennes : machines-outils, véhicules électriques, entraînements directs, alternateurs automobiles ;
  \item fortes puissances : traction ferroviaire, propulsion navale et production d'électricité.
\end{itemize}

L'absence de glissement permet une relation directe entre fréquence électrique et vitesse mécanique. Les aimants permanents offrent une forte densité de puissance ; l'excitation bobinée permet d'agir sur le flux et la puissance réactive.

\subsection{Modèle simplifié d'une phase}

Le modèle monophasé équivalent comporte une f.é.m. interne $\underline E$, la résistance statorique $R_s$ et la réactance synchrone $X_s=L_s\omega$.

\begin{center}\TLSchemaEquivalentMS\end{center}

En convention récepteur, pour un moteur :
\[
\boxed{\underline V=\underline E+(R_s+jX_s)\underline I.}
\]
En négligeant la résistance :
\[
\boxed{\underline V=\underline E+jX_s\underline I.}
\]

La f.é.m. est proportionnelle au flux et à la vitesse :
\[
\boxed{E=K_e\Omega_s=p\Phi\Omega_s.}
\]

\begin{TLMachineWarning}{Domaine de validité}
Le modèle à inductance synchrone unique est établi en régime permanent et convient surtout aux machines à pôles lisses. Les machines saillantes demandent généralement deux réactances, directe et transverse.
\end{TLMachineWarning}

\subsection{Diagramme de Behn--Eschenburg}

Le diagramme de Behn--Eschenburg est le diagramme de Fresnel associé à la loi des mailles.

\begin{center}\TLDiagrammeBehnEschenburg\end{center}

Les angles utilisés sont :
\begin{itemize}
  \item $\varphi$ : déphasage entre $\underline I$ et $\underline V$ ;
  \item $\delta$ : angle interne entre $\underline E$ et $\underline V$ ;
  \item $\psi$ : déphasage entre $\underline I$ et $\underline E$.
\end{itemize}

Lorsque $R_s$ est négligée, le vecteur $jX_s\underline I$ est orthogonal au courant. La valeur de $E$, réglée par l'excitation pour une machine bobinée, modifie le facteur de puissance :
\begin{itemize}
  \item machine sous-excitée : comportement généralement inductif ;
  \item machine surexcitée : comportement capacitif possible ;
  \item excitation adaptée : facteur de puissance proche de l'unité.
\end{itemize}

\begin{TLMachineMethod}{Construire un diagramme de Behn--Eschenburg}
\begin{enumerate}
  \item Choisir la grandeur de référence, souvent $\underline E$ ou $\underline V$.
  \item Placer le courant avec le déphasage imposé.
  \item Construire $R_s\underline I$, colinéaire au courant.
  \item Construire $jX_s\underline I$, en quadrature avance.
  \item Fermer le polygone vectoriel avec la loi des mailles.
  \item Lire les angles $\varphi$, $\psi$ et $\delta$ ainsi que les modules recherchés.
\end{enumerate}
\end{TLMachineMethod}

\subsection{Puissance et couple électromagnétiques}

La puissance électrique active absorbée par une machine triphasée équilibrée est :
\[
P=3VI\cos\varphi.
\]
La puissance électromagnétique peut également s'écrire :
\[
\boxed{P_{em}=3EI\cos\psi.}
\]
On en déduit :
\[
\boxed{C_{em}=\frac{P_{em}}{\Omega_s}=3K_eI\cos\psi.}
\]

Si la résistance statorique est négligée et pour une machine non saillante :
\[
\boxed{P_{em}=3\frac{VE}{X_s}\sin\delta,}
\qquad
\boxed{C_{em}=\frac{3VE}{X_s\Omega_s}\sin\delta.}
\]

Le couple est maximal pour $|\delta|=\pi/2$. Au-delà, la machine entre dans une zone instable et risque le \textbf{décrochage}.

L'angle géométrique $\theta$ entre le champ rotorique et le champ statorique est complémentaire de $\psi$ :
\[
\theta+\psi=\frac{\pi}{2},
\qquad
\sin\theta=\cos\psi.
\]
Ainsi :
\[
C_{em}=3K_eI\sin\theta.
\]

\begin{TLMachineWarning}{Décrochage}
Une machine synchrone directement alimentée par un réseau ne peut pas accepter un couple résistant supérieur à son couple maximal. Le rotor décroche du champ tournant et la machine doit être arrêtée puis redémarrée.
\end{TLMachineWarning}

\begin{TLMachineExercise}{Identification d'une MS à aimants permanents}
Une MS possède huit pôles, des enroulements couplés en étoile et un courant nominal $I_N=155\ \mathrm A$. Les essais donnent :
\begin{itemize}
  \item résistance mesurée entre deux phases : $r=0{,}06\ \Omega$ ;
  \item à vide, à $1500\ \mathrm{tr\,min^{-1}}$ : $E=37\ \mathrm V$ par phase ;
  \item en moteur, à $1500\ \mathrm{tr\,min^{-1}}$ : $I=185\ \mathrm A$, $V=49\ \mathrm V$ et $\psi=0$.
\end{itemize}
Déterminer $f$, la résistance d'une phase, la constante $K_e$, l'inductance synchrone puis construire le diagramme vectoriel pour un fonctionnement à $5000\ \mathrm{tr\,min^{-1}}$.
\end{TLMachineExercise}

\subsection{Bilan de puissance}

Les pertes d'une machine synchrone sont :
\begin{itemize}
  \item pertes Joule statoriques : $P_{Js}=3R_sI^2$ ;
  \item pertes fer $P_{fe}$ ;
  \item pertes mécaniques $P_m$ ;
  \item pertes d'excitation $P_e=U_eI_e=R_eI_e^2$ lorsque le rotor est bobiné.
\end{itemize}

\begin{center}\TLBilanPuissanceMS\end{center}

En fonctionnement moteur :
\[
P_{em}=P_a-P_{Js}-P_{fe},
\qquad
P_u=P_{em}-P_m,
\qquad
\boxed{\eta=\frac{P_u}{P_a+P_e}.}
\]
Pour une machine à aimants permanents, $P_e=0$.

\subsection{Variation de vitesse}

La vitesse étant imposée par la fréquence, la variation de vitesse exige une alimentation à fréquence variable. Le contrôle peut agir sur :
\begin{itemize}
  \item le flux ou l'excitation, si le rotor est bobiné ;
  \item l'amplitude du courant statorique ;
  \item l'angle $\psi$ entre courant et f.é.m.
\end{itemize}

Pour une machine à aimants, l'objectif usuel est de maximiser le couple par ampère en imposant $\psi=0$, donc $\cos\psi=1$.

\subsection{Moteur brushless et autopilotage}

Une machine synchrone autopilotée associe :
\begin{itemize}
  \item un onduleur triphasé ;
  \item une mesure ou une estimation de la position rotorique ;
  \item une commande qui synchronise les courants statoriques avec le rotor.
\end{itemize}

\begin{center}\TLChaineBrushless\end{center}

La commande impose une référence d'amplitude $I_{ref}$ et un angle d'autopilotage $\psi_{ref}$. La fonction assurée par l'électronique est comparable à celle du collecteur mécanique d'une machine à courant continu, d'où le nom \textbf{brushless}.

\begin{TLMachineExercise}{Unité de vissage brushless}
Une MS possède $p=3$, $L=7{,}5\ \mathrm{mH}$ et $K_e=0{,}13\ \mathrm{V/(rad\,s^{-1})}$. Elle tourne à $3600\ \mathrm{tr\,min^{-1}}$.
\begin{enumerate}
  \item Calculer $\omega$ et $f$ des grandeurs statoriques.
  \item Montrer que $C=K_cI\cos\psi$ et déterminer $K_c$.
  \item Choisir $\psi$ pour minimiser le courant en mode moteur puis en freinage.
  \item Tracer les diagrammes de Behn--Eschenburg correspondants.
  \item Pour $N=3000\ \mathrm{tr\,min^{-1}}$ et $I=6{,}7\ \mathrm A$, déterminer le facteur de puissance, la tension simple et le couple.
\end{enumerate}
\end{TLMachineExercise}

\section{Synthèse et démarche d'étude}
\label{sec:masms-6}

\begin{TLMachineMethod}{Analyser une machine alternative}
\begin{enumerate}
  \item Lire la plaque signalétique et identifier le couplage.
  \item Déterminer $p$, $\Omega_s$ et les grandeurs triphasées.
  \item Pour une MAS, calculer le glissement et exploiter le schéma équivalent.
  \item Pour une MS, établir le diagramme de Behn--Eschenburg.
  \item Construire le bilan de puissance et calculer le rendement.
  \item Déterminer le couple et le point de fonctionnement avec la charge.
  \item Vérifier les contraintes thermiques et le domaine stable.
  \item Choisir la stratégie de variation de vitesse : $V/f$, commande vectorielle ou autopilotage.
\end{enumerate}
\end{TLMachineMethod}

\begin{center}
\small
\begin{tabularx}{\linewidth}{>{\bfseries}p{3.2cm}X X}
\toprule
Critère & Machine asynchrone & Machine synchrone \\
\midrule
Vitesse & différente de $\Omega_s$ ; dépend du glissement & égale à $\Omega_s$ en régime établi\\
Rotor & cage ou rotor bobiné & aimants ou excitation bobinée\\
Couple & créé par courants rotoriques induits & interaction entre champs statorique et rotorique\\
Commande simple & robuste, loi $V/f$ & nécessite synchronisation avec le rotor\\
Risque principal & échauffement, zone instable & décrochage\\
Atouts & simplicité, coût, robustesse & rendement, précision de vitesse, densité de puissance\\
\bottomrule
\end{tabularx}
\end{center}

\section{Sources}
\label{sec:masms-7}

Ce cours a été élaboré à partir du document pédagogique original et de ressources de collègues de l'UPSTI. Les exercices conservent les références de concours indiquées dans le document source.
